\section{Proof of Theorem~\ref{thm-der-adaptive}} \label{app-derand}

Throughout, $d = \dig(K,m)$ denotes the digest of $m$ under the pre-hash key $K$, and the slot of a query on $m$ is $d$. Fix a PPT adversary $\advA$ making at most $q_s$ signing queries, and write $N \le q_s$ for the number it makes in a given execution. Signing is deterministic, so a repeated query is answered by replaying the earlier answer in every game below, and the hybrid runs over the distinct queried messages.

\heading{From the real game to the lossy chain.} The real game is the EUF-CMA game for $\SchDScheme[\GG,\HF^{\mathsf{dio}}]$, in which the pre-hash key is generated in injective mode. We add one abort. The event $\mathsf{Bad}_{\mathsf{col}}$ occurs if two distinct messages among the signing queries and the forgery message have equal digests; when it occurs the game outputs 0. In injective mode distinct messages have distinct digests unless key sampling fails, so the abort costs at most $\nu_{\mathsf{inj}}(\secp,r)$.

There is no repeated-slot abort here. In Section~\ref{sec-adaptive} the slot is a message--salt pair and two queries can collide on it; here the slot is the digest, two queries with the same digest have the same message after the collision abort, and such a query is answered by replaying. This is the term the salt costs and derandomizing saves.

Next switch the pre-hash key to lossy mode at the same $r$. The distinguisher receives $K$, samples $x$, $k_0$, and the four PRF keys itself, builds the obfuscated circuit, runs the game including the abort, and outputs its result; it is PPT, so the switch costs $\AdvTLF{\TLF,r}{\advD_{\mathsf{mode}}}$. In every game from here on, an execution that reaches $\Finalize$ without aborting has distinct digests for distinct relevant messages.

\heading{Uniform nonces.} Replace each honest nonce $r_m = \PRFEv_0(k_0,m)$ by an independent uniform value of $\Zp$, once and for all. The key $k_0$ appears nowhere else, so a single distinguisher bounds the change by $\AdvPRF{\PRF_0}{\advD_0}$. Honest nonce points are now uniform, as they are for a randomized signer.

\heading{The outer games.} For $0 \le i \le q_s$, let $\Gm_i$ answer the first $i$ distinct queried messages by the trigger procedure, returning $(W_{d}, s_{d})$ for $d$ the digest of the message, and the remaining ones honestly. Every game publishes the honest circuit of Construction~\ref{constr-der}, retyped to take the digest as its second input; hardwired circuits occur only inside the reductions between adjacent games. Game $\Gm_0$ is the lossy game with the abort, and $\Gm_{q_s}$ answers every query without the secret exponent.

\heading{Changing query $j$.} Fix $j$ and restrict to executions with $N \ge j$. At each endpoint, introduce the event
$$
\mathsf{Bad}_{\mathsf{eq}} := [R_j = W_{d_j}],
\qquad
R_j := g^{r_j},
\qquad
W_{d_j} := g^{s_{d_j}}X^{-c_{d_j}},
$$
and output 0 when it occurs. Conditioned on the view before query $j$, the point $R_j$ is uniform and independent of $W_{d_j}$, so the two endpoint aborts cost at most $2/p$. Every intermediate game applies the same rule to its current values, and when a PRF challenge changes $c_{d_j}$ or $s_{d_j}$ the game recomputes $W_{d_j}$ and the event; this is part of that hybrid.

\heading{The digest guess.} The reduction guesses the digest of query $j$. At key generation it computes $S \getsr \TLFEnum(\mathsl{td})$, removes duplicates, encodes the entries, and extends the list to length exactly $r$ by repeating a fixed string whose first bit is $1$, which is never an encoding of a range element. Call the result $L_S$. It chooses $\widehat\imath \getsr [r]$, sets $\widehat d \gets L_S[\widehat\imath\,]$, and aborts at the end unless $d_j = \widehat d$. Except with probability $\nu_{\mathsf{enum}}(\secp,r)$ the list $S$ contains every digest the lossy key can produce, so $d_j$ occurs in $L_S$, at exactly one position; the index is uniform and independent of everything else, so on those executions the guess is correct with probability exactly $1/r$. Writing $E_b := [N \ge j \land \Gm_b \Rightarrow 1]$ for $b \in \{j-1,j\}$,
$$
\Pr[E_b] \;\le\; r \cdot \Pr[\widehat d = d_j \land E_b] + \nu_{\mathsf{enum}}(\secp,r) \;.
$$
The guessed games fix the slot $\widehat d$ before key generation and output 0 when the guess is wrong, so on the executions with a correct guess they agree with $\Gm_b$. A bound of $\epsilon_j$ on the difference of the guessed games therefore gives
$$
\big|\Pr[E_{j-1}] - \Pr[E_j]\big| \;\le\; r\,\epsilon_j + 2\nu_{\mathsf{enum}}(\secp,r) \;.
$$
This is the only guessing in the proof, and it is where the pre-hash earns its place: without it the guessed value would be the message, at a cost of $2^{\msglen}$.

\begin{lemma} \label{lem-swap}
Let $T_0=(R_0,v_0,s_0)$ and $T_1=(R_1,v_1,s_1)$ be independent and identically distributed, with $v_b,s_b\getsr\Zp$ and $R_b=g^{s_b}X^{-v_b}$, conditioned on $R_0\ne R_1$. Let $V$ be any random variable that depends on $(T_0,T_1)$ only through the unordered set $\{(R_0,v_0),(R_1,v_1)\}$ and through data whose distribution is unchanged when $T_0$ and $T_1$ are exchanged. Then $(V,R_0,s_0)$ and $(V,R_1,s_1)$ are identically distributed.
\end{lemma}

\begin{proof}
Sample the unordered pair $\{T_0,T_1\}$ together with an independent uniform bit $\beta$, and let $T_0,T_1$ be its two elements in the order given by $\beta$. This reproduces the stated distribution, since the two triples are independent, identically distributed, and distinct in their first components. By hypothesis $V$ is a function of the unordered pair and of swap-invariant data, so $V$ is independent of $\beta$. Hence $(V,R_\beta,s_\beta)$ and $(V,R_{1-\beta},s_{1-\beta})$ have the same distribution, which is the claim.
\end{proof}

\heading{The hybrid chain for query $j$.} Set
$$
t = \enczero{\widehat d},
\qquad
R = g^{r_j},
\qquad
t_R = \encone{R,\widehat d},
\qquad
W = g^{s_1}X^{-c},
\qquad
t_W = \encone{W,\widehat d},
$$
where $c = \PRFEv_3(k_3,t)$ and $s_1 = \PRFEv_4(k_4,t)$, and let the main-branch values at $t_R$ be
$$
a = \PRFEv_1(k_1,t_R),
\quad
b = \PRFEv_2(k_2,t_R),
\quad
v = \big(f(RX^{a}g^{b})+a\big)\bmod p,
\quad
s_0 = (r_j + vx)\bmod p .
$$
The two candidate records are $(R,v,s_0)$ and $(W,c,s_1)$. The chain runs as follows. Replace the honest circuit by the two-point circuit that hardwires $\{(t_R,v),(t_W,c)\}$ and punctures $k_1,k_2$ at $T = \{t_R,t_W\}$ and $k_3,k_4$ at $t$; make the two $\PRF_1$ values on $T$ uniform, then the two $\PRF_2$ values; replace $v$ by a uniform value and discard $a,b$, which is exact because $b$ makes $RX^{a}g^{b}$ uniform and independent of $a$, and adding $a$ makes $v$ uniform and independent of $R$; change variables so that the first record is generated by sampling $v,s_0$ and setting $R = g^{s_0}X^{-v}$; make $\PRFEv_3(k_3,t)$ uniform, then $\PRFEv_4(k_4,t)$; exchange the two records; and reverse every step with the second record in place of the first. Each direction invokes $\iO$ once and each of the four puncturable PRFs once, so
$$
\epsilon_j
=
\sum_{\beta\in\{0,1\}}
\Big(
\AdvIO{}{\advD^{\mathsf{io}}_{j,\beta}}
+\sum_{i=1}^{4}\AdvPRF{\PRF_i}{\advD^{\mathsf{prf}}_{j,\beta,i}}
\Big)
+\frac{2}{p} \;.
$$
The circuits compared at the two iO steps are functionally equal: on input $(W,\widehat d)$ both return $c$; on $(R,\widehat d)$ the trigger comparison fails because $\mathsf{Bad}_{\mathsf{eq}}$ has been excluded, and both return $v$; at every other input with second component $\widehat d$ the comparison fails and both use the main branch; and elsewhere functionality preservation of puncturing applies. The reductions holding $k_3\{t\},k_4\{t\}$ must trigger-sign at earlier slots, and a query before $j$ whose slot equals $\widehat d$ cannot be answered with the punctured key. It does not need to be: such a query has the same digest as query $j$, so either the guess is wrong or the collision abort has fired, and the output is 0 in both cases, so the reduction outputs 0 as soon as this occurs.

The exchange step is Lemma~\ref{lem-swap} applied with $V$ the adversary's view apart from the response to query $j$. The retained state is of the kind the lemma allows: the temporary circuit holds the canonically ordered set of two point--value pairs and nothing that distinguishes them, the punctured keys and $t$ are independent of the pair, and no separate copy of $s_0$ or $s_1$ survives, since $r_j$, $a$, $b$, and the unused challenge values are discarded before the comparison.

Summing over the queries,
$$
\big|\Pr[\Gm_0 \Rightarrow 1] - \Pr[\Gm_{q_s} \Rightarrow 1]\big|
\;\le\;
r\sum_{j=1}^{q_s}\epsilon_j + 2q_s\,\nu_{\mathsf{enum}}(\secp,r) \;.
$$

\heading{A main-branch forgery.} Let $E_{\mathsf{main}}$ be the event that the forgery is valid and its hash value comes from the main branch. Adversary $\advB_{\mathsf{main}}$ receives a CDL challenge $h$, sets $X \gets h$, samples the pre-hash key in lossy mode and the PRF keys, publishes the obfuscated circuit, and answers every signing query by the trigger procedure, which uses no secret exponent. Given a valid forgery $(R^*,s^*)$ on $m^*$ it computes $a^*,b^*$ from $k_1,k_2$ at $\encone{R^*,\dig(m^*)}$ and returns $\big(R^*X^{a^*}g^{b^*},\,(s^*+b^*)\bmod p\big)$, which wins the CDL game by the computation of Section~\ref{sec-io-constr}; the zero-free $f$ ensures the solution is accepted. Hence $\Pr[E_{\mathsf{main}}] \le \AdvCDL{\GG,g,f}{\advB_{\mathsf{main}}}$.

\heading{A trigger-branch forgery.} Let $E_{\mathsf{trig}}$ be the event that the forgery is valid and takes the trigger branch, so that $R^* = W_{d^*}$ for $d^* = \dig(m^*)$. The reduction builds $L_S$ as above, chooses $\widehat\imath \getsr [r]$, sets $\widehat d \gets L_S[\widehat\imath\,]$ before publishing the verification key, and aborts unless $\widehat d = d^*$.

The order of the steps matters. In $\Gm_{q_s}$ the published circuit is the honest one and does not depend on $\widehat d$, so the guess is independent of the adversary's view, and except with probability $\nu_{\mathsf{enum}}(\secp,r)$,
$$
\Pr[E_{\mathsf{trig}} \land \widehat d = d^*] = \frac{1}{r}\Pr[E_{\mathsf{trig}}] \;.
$$
We record this relation before changing the circuit. Only then does the reduction puncture $k_4$ at $\widehat t := \enczero{\widehat d}$ and publish the circuit that, on inputs whose slot is $\widehat d$, compares against a hardwired point $Z X^{-c}$ with $c = \PRFEv_3(k_3,\widehat t)$, and otherwise behaves as before. With $Z = g^{\widehat s}$ for $\widehat s = \PRFEv_4(k_4,\widehat t)$ this circuit computes the same function as the honest one, so iO justifies the switch, and puncturable-PRF security at $\widehat t$ replaces $\widehat s$ by a uniform value, after which $Z$ may be set to a DL challenge $Y$, which is uniform in $\GG$. Signing never needs slot $\widehat d$: when the guess is correct, a signing query at that slot would have digest $d^*$, hence either a distinct message with the same digest, which the collision abort excludes, or $m^*$ itself, which freshness excludes. If the forgery takes the trigger branch at the guessed slot then $g^{s^*} = \big(YX^{-c}\big)X^{c} = Y$, so the reduction returns $s^*$. Hence
$$
\Pr[E_{\mathsf{trig}}]
\le
r\Big(
\AdvDL{\GG,g}{\advB_{\mathsf{trig}}}
+\AdvIO{}{\advD^{\mathsf{io}}_{\mathsf{trig}}}
+\AdvPRF{\PRF_4}{\advD^{\mathsf{prf}}_{\mathsf{trig}}}
\Big)
+\nu_{\mathsf{enum}}(\secp,r) \;.
$$

\heading{Conclusion.} A valid forgery takes one of the two branches. Collecting the terms,
$$
\begin{aligned}
&\AdvEUFCMA{\SchDScheme[\GG,\HF^{\mathsf{dio}}]}{\advA}\\
&\quad\le
\AdvTLF{\TLF,r}{\advD_{\mathsf{mode}}}
+\AdvPRF{\PRF_0}{\advD_0}
+\AdvCDL{\GG,g,f}{\advB_{\mathsf{main}}}\\
&\qquad+r\Big(
\AdvDL{\GG,g}{\advB_{\mathsf{trig}}}
+\AdvIO{}{\advD^{\mathsf{io}}_{\mathsf{trig}}}
+\AdvPRF{\PRF_4}{\advD^{\mathsf{prf}}_{\mathsf{trig}}}
\Big)\\
&\qquad+r\sum_{j=1}^{q_s}\epsilon_j
+\nu_{\mathsf{inj}}(\secp,r)
+(2q_s+1)\,\nu_{\mathsf{enum}}(\secp,r) \;,
\end{aligned}
$$
which is the bound of the theorem. Every reduction runs in time polynomial in $\secp$ and $r$, the enumeration of the lossy image being the only super-polynomial cost. \qed

\paragraph*{Proof of Corollary~\ref{cor-der-adaptive}.} We bound the display under the hypotheses of the corollary. Each CDL, iO, and puncturable-PRF advantage is at most $\eps$, and so is the DL advantage, by the tight reduction from DL to zero-free CDL of Section~\ref{sec-derand-adaptive}. Each $\epsilon_j$ contains two iO and eight puncturable-PRF advantages together with $2/p$, and the trivial discrete-logarithm algorithm gives $1/p \le \eps$, so $\epsilon_j \le 12\eps$. Since $r$ is super-polynomial and $q_s$ polynomial, $q_s \le r$ for large $\secp$, and the display is at most
$$
\AdvTLF{\TLF,r}{\advD_{\mathsf{mode}}}
+\AdvPRF{\PRF_0}{\advD_0}
+\big(1 + 3r + 12r^2\big)\eps
+\nu_{\mathsf{inj}}(\secp,r)
+(2q_s+1)\,\nu_{\mathsf{enum}}(\secp,r) \;.
$$
The first two terms are negligible because $\TLF$ and $\PRF_0$ are secure and the two reductions are PPT. The third is at most $16 r^2 \eps$, which is negligible by hypothesis. The remaining terms are negligible at $r$. \qed
